\documentclass[11pt]{article}
\usepackage[a4paper,margin=1.15in]{geometry}
\usepackage{amsmath,amssymb}
\usepackage{enumitem}
\usepackage{array}
\usepackage{xcolor}
\usepackage{mdframed}
\usepackage{hyperref}
\hypersetup{hidelinks}
\newcommand{\st}[1]{\textsf{\small[#1]}}
\title{Pointless Coherence: Order Without a Centre}
\author{Patrick S.\ Mahon}
\date{October 2026 \quad (draft)}
\begin{document}
\maketitle

\begin{mdframed}[backgroundcolor=black!4,linewidth=0.4pt]
\textbf{The question.} When do local views, each correct as far as it goes, add
up to one world? And when they do not, what are they still telling us?

\smallskip
Part~I gives the mathematical and philosophical core. Part~II asks where the
question is lived: in ecosystems with no species on top, and in societies with
no one in charge. Part~II is a lens, not an application, and
Section~\ref{sec:distance} says plainly how far the mathematics reaches.
Status labels: \st{Lean} machine-checked, \st{hand} proved by hand, otherwise
conjecture.
\end{mdframed}

\part*{Part I. The essence}

\section{The stance}
Answers come first. Questions, and patterns of answers to them that are
consistent, are taken as primitive. A \emph{world} is something answers have to
\emph{earn} by fitting together. It may be an actual point, or an ideal point at
infinity. Nothing is assumed to stand behind the answers except the bare space of
possibilities.

The patterns studied here are \emph{complete} (every question that can be asked
gets a yes or a no) and \emph{coherent} (answers to disjoint questions add up).
In ordinary logic, complete and coherent always means there is a world: the
completeness theorem~\cite{godel1930}. That relies on every ``and'' being
askable. The theory studies what happens when it is not, and so it studies one
gap:
\begin{center}\emph{complete and coherent, yet about no world.}\end{center}

\section{The picture}
Penrose's staircase~\cite{penrose1991}. Each flight of stairs is a \emph{window}:
a view of two landings that says ``this one is higher''. Every flight goes up,
and the flights close into a loop. Each view is fine on its own, but no building
satisfies all of them. The board lets you ask questions about one flight at a
time, and anything built from those by ``not'' and by joining questions that
never both say yes. Joint questions across flights are not automatically
askable.

\section{What is known}
\begin{center}\small
\begin{tabular}{|>{\raggedright\arraybackslash}p{6.4cm}|>{\raggedright\arraybackslash}p{4.4cm}|>{\raggedright\arraybackslash}p{2.6cm}|}
\hline
\textbf{regime} & \textbf{coherent answers earn a world?} & \textbf{status}\\\hline
finitely many possible heights & yes, an actual one & \st{hand}\\
infinitely many, answers add over the whole infinity & yes, an actual one & \st{hand; Lean on rings}\\
infinitely many, finite adding, no loop & yes, an ideal one & \st{hand}\\
infinitely many, finite adding, a loop & \textbf{no}: pointless answers exist & \st{Lean + hand}\\\hline
\end{tabular}
\end{center}

\noindent\textbf{Theorem.} \emph{Coherent answers can be about no world exactly
when the views form a loop, there are infinitely many possibilities, and answers
add up only finitely.}

\medskip
\noindent In Aristotle's terms~\cite{aristotle}, finite adding is the
\emph{potential} infinite. Demanding the actual, completed infinity forces a world
back into existence. ``Ideal'' and ``pointless'' match the standard notions of
coherent and non-coherent finitely additive states~\cite{bhaskararao1983}. For
ordinary probabilities and finitely many outcomes, the tree/loop line is
classical~\cite{vorobev1962,bfmy1983}.

\section{Why: what the board contributes}
On the triangle (three landings, three flights) the contradiction could only be
exposed by one question: \emph{do the three heights go round the staircase's
way?} That question can never be built, not even in infinitely many steps
\st{Lean}. Two facts do the work.
\begin{itemize}[nosep]
  \item \emph{Balance.} Pick two candidate heights for each landing. That gives
    the eight corners of a box. Every buildable question says yes equally often
    on the even and the odd corners.
  \item \emph{The lift.} Raising the lowest landing above the highest never
    changes the answer to a buildable question, because the box joining the two
    configurations has its other six corners going the other way round.
\end{itemize}
So the coherence survives because the one question that would break it is
unaskable. Beyond that one exclusion, whether answers are pointless depends only
on whether the windows' local answers fit together. In every case checked, any
local answers that agree where they overlap extend to the whole board
\st{hand; Lean for the triangle cases}. The nearest precedent is
Drav\'eck\'y and \v{S}ipo\v{s}'s example~\cite{draveckysipos1980} and the work
of Navara and Pt\'ak~\cite{navaraptak1983,navara1984}. There, answers that
add over the whole infinity can be pointless on a sub-board, but they do not
extend to the full board.

\section{What creates impossibility, and what does not}
\begin{itemize}[nosep]
  \item \emph{Renaming never does.} A window that only says ``this landing is
    that one, renamed'' can be contracted away \st{hand, via Katětov's
    lemma~\cite{katetov1967}}.
  \item \emph{Motion never does.} Symmetry and walking the stairs only move
    worlds around. The twist lives in the loop, not in the action \st{hand}.
  \item \emph{``Beyond'' can.} A window that draws one landing \emph{beyond} the
    other carries a direction. A loop of such directions is a task that would
    have to finish before it started.
  \item \emph{And beyonds compose.} A triangle whose directions never go round
    can still fail to close. The combined step from the first landing to the
    last must be the composite of the other two in the algebra of
    ultrafilters~\cite{hindmanstrauss,goldberg2025}, and one window can ask for a
    different composite \st{hand, not yet independently checked}. So the true
    invariant is a \emph{composition law}, with the directed cycle as its
    simplest shadow. Its exact form is the main open question.
\end{itemize}

\section{The lineage}
\begin{itemize}[nosep]
  \item \emph{Specker}~\cite{specker1960}: a seer's three boxes, any two of
    which can be opened consistently but not all three. That is the triangle,
    and the seed of the Kochen--Specker theorem~\cite{kochenspecker1967}. The
    sheaf-theoretic account of contextuality~\cite{abramskybrandenburger2011}
    and the classification of cycle scenarios~\cite{kharoofipekokay2024} own the
    finite, ordinary-probability side.
  \item \emph{Penrose}~\cite{penrose1991}: impossible figures as a property of
    loops. Here the loop is necessary, and the obstruction is the way steps
    compose.
  \item \emph{Kant}~\cite{kant}: the completed totality of conditions is not
    \emph{given} but \emph{set as a task}. An ideal point is that task, and a
    loop is where the task becomes circular.
\end{itemize}

\part*{Part II. The lived}

\section{The lizards}
On the hills of central California, male side-blotched lizards come in three
throat colours, each with its own strategy. Orange males hold large territories
and take mates from blue males. Blue males guard their mates closely and so beat
the yellow males. Yellow males sneak into orange territories and beat the
oranges. Orange beats blue, blue beats yellow, yellow beats orange. The
frequencies cycle over the years, and no colour wins~\cite{sinervolively1996}.

That is Penrose's staircase, alive: every flight goes up, the stairs close into
a loop, and there is no top. Similar cycles are documented among
toxin-producing strains of \emph{E.\ coli}, where the cycle keeps all three
strains alive when interactions are local~\cite{kerr2002}, and among encrusting
organisms on coral reefs~\cite{bussjackson1979}. Network models proposed that
such cycles can sustain many more species than a ranking
could~\cite{allesinalevine2011}. A later replication found the effect much
weaker~\cite{larsenhaine2015}, so the strong claim is contested.

Ordinary logic says a cycle of ``beats'' is a contradiction, because no ranking
satisfies it. Ecology shows it can be an order. It is just not a
\emph{ranking}. That is the shape of the central finding of Part~I: coherent
answers with no single world behind them.

\section{Ecology: order from local relations}
\paragraph{Cycles of dominance.} An ecosystem does not consult a global
ranking. Each encounter is local: this lizard, that territory, this season.
When local relations form a \emph{tree}, they always add up to one global order.
When they form a \emph{loop}, they can be coherent with no global order at all.

\paragraph{Steps across scales.} Holling's \emph{panarchy}~\cite{panarchy2002}
describes nested adaptive cycles at different speeds: a leaf, a stand, a
landscape. The finding that beyonds compose rhymes with this. The question is not
only whether influences go round in a loop, but whether they \emph{compose}
across scales. It is a rhyme, not yet a theorem about ecosystems.

\paragraph{Partial observation.} Ecological knowledge is assembled from plots,
transects, traps and seasons, each a local view. Do my surveys connect in a
loop? If not, local agreement is enough. If so, they can agree pairwise and
still describe no single landscape.

\section{Anarchism: order without a sovereign}
\paragraph{The stance is already anarchist.} The theory takes \emph{answers
first}. There is no presupposed ground and no sovereign ``real world'' that the
answers report on. A single world, or a single centre, is something local
answers must \emph{earn}. This mirrors an anarchist account of order: it is not
handed down from a centre, it is made, or not made, from relations.
Feyerabend's ``epistemological anarchism''~\cite{feyerabend1975} makes the same
refusal for science.

\paragraph{Hierarchies are trees.} Local agreements arranged as a tree always
add up to one global order. A chain of command is a tree, which is why
hierarchies \emph{look} coherent: structurally, they cannot fail to be. The cost
is that coherence is bought by removing every relation that does not route
through the centre.

\paragraph{Federations are loops.} Kropotkin was a field naturalist before he
was an anarchist. In \emph{Mutual Aid}~\cite{kropotkin1902} he argued from
animals, villages, guilds and medieval cities that cooperation built from local
relations is as real a force in evolution as competition. Elinor Ostrom made
the case empirically~\cite{ostrom1990}. Swiss alpine pastures, Japanese village
forests, Spanish irrigation huertas and Philippine irrigation societies
governed shared resources for centuries with no state or market in charge. One
of her design principles is \emph{nested enterprises}: overlapping local centres
of rule-making. Her \emph{polycentric governance}~\cite{ostrom2010} generalises
it. In the theory's terms these are webs with loops. Pointless coherence is,
structurally, \emph{order without a sovereign}.

\paragraph{How commitments chain.} Local agreements can be fine, even pairwise,
yet fail to add up. That happens not because anyone contradicts anyone, but
because the agreement between the first and last parties must match the
composite of those in between.

\section{The four questions, lived}
\begin{center}\small
\begin{tabular}{|>{\raggedright\arraybackslash}p{3.9cm}|>{\raggedright\arraybackslash}p{4.6cm}|>{\raggedright\arraybackslash}p{4.6cm}|}
\hline
\textbf{question} & \textbf{in an ecosystem} & \textbf{in a commons or federation}\\\hline
Is there a loop? & cycles of dominance, nutrient cycles, feedbacks & overlapping councils, reciprocal obligations\\\hline
Is there a joint question nobody can ask? & ``who is best?'' when dominance cycles & ``who is in charge?'' in a polycentric order\\\hline
Is an idealisation hiding a contradiction? & equilibrium, carrying capacity, ``in the long run'' & ``the general will'', ``the market'' as one agent\\\hline
Circle, or composition? & do influences compose across scales? & do commitments chain consistently?\\\hline
\end{tabular}
\end{center}

\section{Honest distance}\label{sec:distance}
\begin{center}\small
\begin{tabular}{|>{\raggedright\arraybackslash}p{4.2cm}|>{\raggedright\arraybackslash}p{9.2cm}|}
\hline
\textbf{link} & \textbf{status}\\\hline
lizards, \emph{E.\ coli}, reefs & \emph{literal} as structure: a directed cycle of pairwise dominance is exactly the staircase's shape\\\hline
hierarchy as tree, federation as loop & \emph{structural}: the forest/loop theorem applied to the shape of relations\\\hline
answers-first as anarchist epistemology & \emph{philosophical}: a shared refusal of a presupposed centre\\\hline
panarchy and composition & \emph{a rhyme}: suggestive, not yet a theorem\\\hline
\end{tabular}
\end{center}
Part~I works with idealised infinities. With finitely many possibilities its own
theorem says every coherent pattern \emph{is} a world. The lizard cycle is
coherent as a \emph{dynamics} rather than as a set of answers, and real commons
are finite. So the lived reading needs a different bridge: a
\emph{probabilistic} or \emph{dynamical} version, in which ``locally consistent
but globally unrealisable'' is real for finite systems. Belief propagation on
loopy networks is one such setting~\cite{wainwrightjordan2008}. Whether the
composition law has a finite, probabilistic shadow is open, and it may be the
bridge that matters most.

\section{What the lived side asks of the mathematics}
\begin{enumerate}[nosep]
  \item \emph{A finite, probabilistic version.} When are local marginals on a
    web of relations consistent yet unrealisable? Does composition, and not only
    cycles, govern it?
  \item \emph{Scale as the source of ``beyond''.} In the theory, ``beyond every
    number'' creates twist. In ecology and society the analogue may be
    \emph{separation of scales}. Can that be made precise?
  \item \emph{Dynamics, revisited.} The theory found that motion only moves
    worlds. The lizard cycle is a \emph{dynamic} order. Can a cycle in time,
    rather than in the questions, carry coherence without a top?
  \item \emph{Design.} Ostrom asked which institutions hold together. Could the
    theory say which webs of local rules can cohere without a centre, and which
    must fail by composition?
\end{enumerate}

\section*{Questions to carry back into the world}
\begin{enumerate}[nosep]
  \item Do my views form a loop?
  \item Is there a joint question nobody can ask?
  \item Is an idealisation hiding a contradiction?
  \item When things fail to add up, is it a circle, or the way steps compose?
\end{enumerate}

\begin{thebibliography}{99}\small
\bibitem{abramskybrandenburger2011} S.~Abramsky and A.~Brandenburger, The
  sheaf-theoretic structure of non-locality and contextuality, \emph{New J.\
  Phys.} 13 (2011) 113036.
\bibitem{allesinalevine2011} S.~Allesina and J.~M.~Levine, A competitive
  network theory of species diversity, \emph{PNAS} 108(14) (2011) 5638--5642.
\bibitem{aristotle} Aristotle, \emph{Physics}, Book III, ch.~6.
\bibitem{bfmy1983} C.~Beeri, R.~Fagin, D.~Maier and M.~Yannakakis, On the
  desirability of acyclic database schemes, \emph{J.\ ACM} 30(3) (1983)
  479--513.
\bibitem{bhaskararao1983} K.~P.~S.~Bhaskara Rao and M.~Bhaskara Rao,
  \emph{Theory of Charges: A Study of Finitely Additive Measures}, Pure and
  Applied Mathematics 109, Academic Press, London, 1983.
\bibitem{bussjackson1979} L.~W.~Buss and J.~B.~C.~Jackson, Competitive networks:
  nontransitive competitive relationships in cryptic coral reef environments,
  \emph{Amer.\ Natur.} 113 (1979) 223--234.
\bibitem{draveckysipos1980} J.~Drav\'eck\'y and J.~\v{S}ipo\v{s}, On the
  additivity of Gudder integral, \emph{Math.\ Slovaca} 30(3) (1980) 299--303.
\bibitem{feyerabend1975} P.~Feyerabend, \emph{Against Method}, New Left Books,
  London, 1975.
\bibitem{godel1930} K.~G\"odel, Die Vollst\"andigkeit der Axiome des logischen
  Funktionenkalk\"uls, \emph{Monatsh.\ Math.\ Phys.} 37 (1930) 349--360.
\bibitem{goldberg2025} G.~Goldberg, On products of ultrafilters, \emph{J.\
  Symbolic Logic}, published online 2025, doi:10.1017/jsl.2025.10097;
  arXiv:2102.04677.
\bibitem{panarchy2002} L.~H.~Gunderson and C.~S.~Holling (eds.),
  \emph{Panarchy: Understanding Transformations in Human and Natural Systems},
  Island Press, Washington, DC, 2002.
\bibitem{hindmanstrauss} N.~Hindman and D.~Strauss, \emph{Algebra in the
  Stone--\v{C}ech Compactification}, de Gruyter, 1998; 2nd ed.\ 2012.
\bibitem{kant} I.~Kant, \emph{Kritik der reinen Vernunft}, 1781/1787;
  Antinomy of Pure Reason, \S\S7--8.
\bibitem{katetov1967} M.~Kat\v{e}tov, A theorem on mappings, \emph{Comment.\
  Math.\ Univ.\ Carolinae} 8(3) (1967) 431--433.
\bibitem{kerr2002} B.~Kerr, M.~A.~Riley, M.~W.~Feldman and B.~J.~M.~Bohannan,
  Local dispersal promotes biodiversity in a real-life game of
  rock--paper--scissors, \emph{Nature} 418 (2002) 171--174.
\bibitem{kharoofipekokay2024} A.~Kharoof, S.~Ipek and C.~Okay, Extremal
  simplicial distributions on cycle scenarios with arbitrary outcomes,
  arXiv:2406.19961.
\bibitem{kochenspecker1967} S.~Kochen and E.~P.~Specker, The problem of hidden
  variables in quantum mechanics, \emph{J.\ Math.\ Mech.} 17 (1967) 59--88.
\bibitem{kropotkin1902} P.~Kropotkin, \emph{Mutual Aid: A Factor of Evolution},
  William Heinemann, London, 1902.
\bibitem{larsenhaine2015} Larsen and Haine, Monte Carlo replication of
  competitive network theory, abstract, Ecological Society of America annual
  meeting, 2015 (conference abstract; not yet read in full).
\bibitem{navara1984} M.~Navara, Two-valued states on a concrete logic and the
  additivity problem, \emph{Math.\ Slovaca} 34(3) (1984) 329--336.
\bibitem{navaraptak1983} M.~Navara and P.~Pt\'ak, Two-valued measures on
  $\sigma$-classes, \emph{\v{C}asopis P\v{e}st.\ Mat.} 108(3) (1983) 225--229.
\bibitem{ostrom1990} E.~Ostrom, \emph{Governing the Commons: The Evolution of
  Institutions for Collective Action}, Cambridge University Press, 1990.
\bibitem{ostrom2010} E.~Ostrom, Beyond markets and states: polycentric
  governance of complex economic systems, \emph{Amer.\ Econ.\ Rev.} 100(3)
  (2010) 641--672.
\bibitem{penrose1991} R.~Penrose, On the cohomology of impossible figures,
  \emph{Structural Topology} 17 (1991) 11--16; reprinted in \emph{Leonardo} 25
  (1992).
\bibitem{sinervolively1996} B.~Sinervo and C.~M.~Lively, The rock--paper--scissors
  game and the evolution of alternative male strategies, \emph{Nature} 380
  (1996) 240--243.
\bibitem{specker1960} E.~Specker, Die Logik nicht gleichzeitig entscheidbarer
  Aussagen, \emph{Dialectica} 14 (1960) 239--246.
\bibitem{vorobev1962} N.~N.~Vorob'ev, Consistent families of measures and their
  extensions, \emph{Theory Probab.\ Appl.} 7(2) (1962) 147--163.
\bibitem{wainwrightjordan2008} M.~J.~Wainwright and M.~I.~Jordan, Graphical
  models, exponential families, and variational inference, \emph{Found.\
  Trends Mach.\ Learn.} 1(1--2) (2008) 1--305.
\end{thebibliography}

\medskip\noindent\small Full proofs and status: \texttt{pointless\_coherence.tex}.
The theory's current shape: \texttt{theory\_shape.md}. The why: \texttt{WHY.md}.
\end{document}
